\section{A repeated-eigenspace counterexample}\label{sec:counterexample}

This section presents a complete rational instance that refutes a natural simplification in the
hyperplane layer: \emph{taking, for each distinct real eigenvalue of the compression, one eigenvector as
the representative direction}, and deciding feasibility from that representative alone. In the instance
the compression has a two-dimensional real eigenspace, and different projective directions inside it
lead to \emph{different} answers: one direction is feasible, the others tested are infeasible. What is
refuted is exactly that reduction; the example does not show that the repeated case is computationally
hard, nor that geometric simplicity of the real spectrum is necessary for polynomial-time decidability
(Section~\ref{sec6:sub:conclusion}). All computations are exact rational arithmetic: infeasibility of a
direction is certified by a divisibility identity over the function field of the fibre, using the
maximal-minor test of \cite{companion1}, and feasibility of the good direction is certified at a single
rational chart point.

\subsection{Notation}\label{sec6:sub:notation}

We use the notation of the preceding sections. The section lives in the hyperplane layer $c=1$, so
$\ker H$ is one-dimensional, say $\ker H=\mathbb{R}v$ with $v\in K_0$, every row of $\mathcal{L}$
annihilates $v$ and $Av$, and with $U_v=\operatorname{span}\{v,Av\}$ the admissible row spaces obey
$V\subseteq U_v^{\perp}$ and $\dim(V\cap K_0)=k-1$. Write
\begin{equation}\label{sec6:eq:defs}
D_v=K_0\cap U_v^{\perp},\qquad \mathcal{Q}_v=U_v^{\perp}/D_v,\qquad s=q-k+1 .
\end{equation}
Since $L_0$ is empty, the admissible directions are exactly the real eigenvector directions of the
compression, and for fixed $[v]$ the fibre is $\mathbb{P}(\mathcal{Q}_v)$, of dimension $r(v)-1$ with
$r(v)=\dim\mathcal{Q}_v=\rho$ for an eigen direction ($Av\in\operatorname{span}v$) and $\rho-1$ for a
non-eigen direction. In this section $s=1$, so every fibre is a projective space and a single
computation over the function field of the fibre covers it entirely: no chart decomposition is needed.

\subsection{The instance}\label{sec6:sub:instance}

Take
\begin{equation}\label{sec6:eq:data}
n=8,\quad p=2,\quad q=4,\qquad
C=\begin{bmatrix}e_1\\ e_2\end{bmatrix},\qquad L_0\ \text{empty},
\end{equation}
and let the state matrix be
\begin{equation}\label{sec6:eq:A}
A=\begin{bmatrix}
0&0&1&0&0&0&0&0\\
0&0&0&1&0&0&0&0\\
1&-2&1&0&-1&1&1&-1\\
-1&2&1&0&2&-2&-1&2\\
0&0&0&0&1&0&0&0\\
0&0&0&0&0&1&0&0\\
0&0&0&0&0&0&2&0\\
0&0&0&0&0&0&0&3
\end{bmatrix}.
\end{equation}
Its lower-right block is $\operatorname{diag}(1,1,2,3)$, its upper-right block is
$\begin{bsmallmatrix}0&0&0&0\\0&0&0&0\\-1&1&1&-1\\2&-2&-1&2\end{bsmallmatrix}$, its lower-left block
vanishes, and its upper-left block has first two rows giving $CA=[e_3;e_4]$.

It follows that
\begin{equation}\label{sec6:eq:K0}
K_0=\operatorname{span}\{e_5,e_6,e_7,e_8\},
\qquad
N=\begin{bmatrix}
0&0&0&0\\0&0&0&0\\0&0&0&0\\0&0&0&0\\
1&0&0&0\\0&1&0&0\\0&0&1&0\\0&0&0&1
\end{bmatrix},
\end{equation}
hence
\begin{equation}\label{sec6:eq:invariants}
\rho=k=4,\qquad \omega=k-1=3,\qquad s=q-k+1=1 ,
\end{equation}
and, in the coordinates of $K_0$,
\begin{equation}\label{sec6:eq:M}
\begin{aligned}
M&=\pi_{K_0}A|_{K_0}=\operatorname{diag}(1,1,2,3),\\
\chi_M(x)&=(x-3)(x-2)(x-1)^2 .
\end{aligned}
\end{equation}
\begin{equation}\label{sec6:eq:A21}
A_{21}=\pi_{K_0^{\perp}}A|_{K_0}=
\begin{bmatrix}
0&0&0&0\\0&0&0&0\\-1&1&1&-1\\2&-2&-1&2\\
0&0&0&0\\0&0&0&0\\0&0&0&0\\0&0&0&0
\end{bmatrix}.
\end{equation}

The following identities hold in exact arithmetic: $[C;CA]=[e_1;e_2;e_3;e_4]$; $N^{\mathsf{T}}N=I_4$ and
$N^{\mathsf{T}}AN=\operatorname{diag}(1,1,2,3)$; $AN=NM+A_{21}$; and $\chi_M(x)=(x-3)(x-2)(x-1)^2$. The
instance is obtained from a larger example by keeping $C$ and the action of $A$ outside $K_0$ and
replacing only the action on $K_0$, which creates a two-dimensional real eigenspace while leaving
$K_0$, $\rho$, $k$ and $\omega$ unchanged.

\subsection{Direction geometry}\label{sec6:sub:direction}

The factor $(x-1)^2$ produces a two-dimensional eigenspace
$E_1=\ker(M-I)=\operatorname{span}\{e_1,e_2\}\subseteq\mathbb{R}^{4}$, in the coordinates of $K_0$
(corresponding to $e_5,e_6$ in the state space); the eigenspaces of the eigenvalues $2$ and $3$ are
one-dimensional. Directions in $E_1$ are written
\begin{equation}\label{sec6:eq:v}
v(c_1,c_2)=c_1e_5+c_2e_6,\qquad (c_1,c_2)\neq(0,0),
\end{equation}
so the direction space is the projective line $\mathbb{P}(E_1)\cong\mathbb{P}^{1}$.

\textbf{The eigen/non-eigen stratification.} For $c=(c_1,c_2,0,0)^{\mathsf{T}}\in E_1$,
\begin{equation}\label{sec6:eq:A21c}
\begin{aligned}
A_{21}c&=(c_2-c_1)\begin{bmatrix}1&-2&0&0&0&0&0&0\end{bmatrix}^{\mathsf{T}},\\
A_{21}c&=0\quad\iff\quad c_1=c_2 .
\end{aligned}
\end{equation}
Thus $E_1$ contains exactly \emph{one} eigen direction, $[1\!:\!1]$, and every other direction of $E_1$
is a non-eigen direction; this is consistent with the general description, since the eigen directions
inside $E_\mu$ form the linear subspace $\ker A_{21}$. For $[1\!:\!1]$ one checks directly that $Av=v$,
so $\dim U_v=1$.

\textbf{Fibres.} With $s=1$ and $\delta=0$ the fibre over a fixed direction is $\mathbb{P}^{2}$ for a
non-eigen direction ($r(v)=\rho-1=3$) and $\mathbb{P}^{3}$ for the eigen direction
($r(v)=\rho=4$): the eigen direction has $\dim U_v=1$ instead of $2$, hence $\dim U_v^{\perp}=7$
instead of $6$, while $\dim D_v=k-1=3$ is unchanged, so its fibre gains exactly one homogeneous
parameter. The two situations are compared in Table~\ref{sec6:tab:dims}.

\begin{table}[t]
\centering
\caption{Dimensions attached to a direction, for $n=8$ and $k=4$. The two columns differ because
$\dim U_v$ differs: an eigen direction places $Av$ inside $\operatorname{span}v$, so $U_v$ loses one
dimension and its annihilator gains one.}
\label{sec6:tab:dims}
\begin{tabular}{@{}lcc@{}}
\toprule
Quantity & Non-eigen & Eigen\\
\midrule
$\dim U_v$ & $2$ & $1$\\
$\dim U_v^{\perp}=n-\dim U_v$ & $6$ & $7$\\
$r(v)=\dim U_v^{\perp}-(k-1)$ & $3$ & $4$\\
Fibre $\mathbb{P}^{r(v)-1}$ & $\mathbb{P}^{2}$ & $\mathbb{P}^{3}$\\
\bottomrule
\end{tabular}
\end{table}

\textbf{Example: the direction $(1\!:\!0)$.} Here $v=e_5$ and
$Av=Ae_5=(0,0,-1,2,1,0,0,0)^{\mathsf{T}}$, so in the coordinates of $K_0$
\begin{equation}\label{sec6:eq:Dv}
\begin{aligned}
D_v&=K_0\cap U_v^{\perp}=\{z\in K_0:\ z\cdot v=z\cdot Av=0\}\\
   &=\operatorname{span}\{e_6,e_7,e_8\},
\end{aligned}
\end{equation}
$\mathcal{Q}_v$ is three-dimensional, the fibre is $\mathbb{P}^{2}$, and the admissible row spaces are
$\row\mathcal{L}=\operatorname{span}\{e_6,e_7,e_8\}\oplus\langle x\rangle$ with
$0\neq x\in\mathcal{Q}_v$. Writing $\mathcal{L}=[\ell_1;\ell_2;\ell_3;x]$ with $\ell_1,\ell_2,\ell_3$ a
basis of $D_v$, we form $H=LN$ and $H_A=LAN$; the maximal minors $m_I(x,\lambda)$ are the
$\binom{4}{3}\binom{4}{3}=16$ minors of order $3$, some identically zero, each a polynomial in
$(x,\lambda)$.

\subsection{Decisions for six directions}\label{sec6:sub:six}

We test six directions of $\mathbb{P}(E_1)$; the third is the eigen direction, the other five are
non-eigen.

\begin{table*}[t]
\centering
\caption{Decisions for six projective directions of the repeated eigenspace $E_1$. The fourth column
counts the $3\times3$ maximal minors that do not vanish identically over the function field; the fifth
column is the $\lambda$-degree of the greatest common divisor of all maximal minors over that field.
Degree $0$ means that no common root exists, hence that the fibre is generically good and, by the
full-rank leading-coefficient lemma of Section~\ref{sec:decision}, good at every point of the chart.}
\label{sec6:tab:dirs}
\begin{tabular}{@{}cccccl@{}}
\toprule
Direction $(c_1\!:\!c_2)$ & Type & Fibre & Nonzero maximal minors & $\deg_\lambda\gcd$ & Verdict\\
\midrule
$(1\!:\!0)$ & non-eigen & $\mathbb{P}^{2}$ & $2$ & $2$ & entire fibre infeasible\\
$(0\!:\!1)$ & non-eigen & $\mathbb{P}^{2}$ & $2$ & $2$ & entire fibre infeasible\\
$(1\!:\!1)$ & \textbf{eigen} & $\mathbb{P}^{3}$ & $4$ & $\mathbf{0}$ & \textbf{a feasible point exists}\\
$(1\!:\!2)$ & non-eigen & $\mathbb{P}^{2}$ & $4$ & $2$ & entire fibre infeasible\\
$(2\!:\!-1)$ & non-eigen & $\mathbb{P}^{2}$ & $4$ & $2$ & entire fibre infeasible\\
$(3\!:\!1)$ & non-eigen & $\mathbb{P}^{2}$ & $4$ & $2$ & entire fibre infeasible\\
\bottomrule
\end{tabular}
\end{table*}

\textbf{Observation.} Inside one and the same two-dimensional eigenspace, the fibre over $(1\!:\!1)$
contains a feasible row space, whereas every fibre over the other five directions is infeasible, so
feasibility on $\mathbb{P}(E_1)$ is \emph{not constant}. The differing numbers of nonzero minors ($2$
versus $4$) are a consequence of the particular directions and do not affect the decision, which uses
the gcd of \emph{all} minors.

\textbf{Why the eigen direction is the good one.} The two phenomena are related but not equivalent: the
eigen direction is the only one whose fibre is a $\mathbb{P}^{3}$ instead of a $\mathbb{P}^{2}$, so its
parameter space is one dimension larger and has more room to avoid the bad locus. This instance
therefore also illustrates the general statement that the eigen and non-eigen strata must be handled
separately; it does not assert that an eigen direction is always the feasible one.

\subsubsection{Location of the common roots}\label{sec6:sub:roots}

The common factors of the infeasible directions have the $\lambda$-parts of Table~\ref{sec6:tab:roots}:
for $(1\!:\!0)$ the polynomial $(\lambda-1)(\lambda-3)$, and a power of $(\lambda-1)$ for the other
directions, so all common roots lie in the spectrum $\{1,2,3\}$ of the compression.

\begin{table}[t]
\centering
\caption{Location of the common roots. All of them lie in the spectrum $\{1,2,3\}$ of the compression.}
\label{sec6:tab:roots}
\begin{tabular}{@{}lll@{}}
\toprule
Direction & $\lambda$-part of the common factor & Common roots\\
\midrule
$(1\!:\!0)$ & $(\lambda-1)(\lambda-3)$ & $\lambda=1,\ 3$\\
direction of $\mu=2$ & $(\lambda-1)$ & $\lambda=1$\\
direction of $\mu=3$ & $(\lambda-1)^{2}$ & $\lambda=1$\\
\bottomrule
\end{tabular}
\end{table}

\textbf{Range of the computation.} The direction space $\mathbb{P}(E_1)$ is one-dimensional and we
have taken six of its points, including the unique eigen direction; this establishes non-constancy but
is \emph{not} a classification of all directions of $\mathbb{P}(E_1)$.

\subsection{Certificates}\label{sec6:sub:cert}

\textbf{Infeasible directions.} For each of the five non-eigen directions of Table~\ref{sec6:tab:dirs},
over the function field $\mathbb{Q}(x_1,\dots,x_{r(v)-1})$ the exact computation gives
$\deg_\lambda g=2$ and $g\mid m_I$ for every maximal minor $I$, divisibility being verified by division
with remainder over that field (the remainder is zero). For $(1\!:\!0)$, with homogeneous coordinates
$(x_0\!:\!x_1\!:\!x_2)$ on the fibre, the gcd is exactly
\begin{equation}\label{sec6:eq:gexplicit}
g=x_2\,(\lambda-1)(\lambda-3).
\end{equation}
Its $\lambda$-part is $(\lambda-1)(\lambda-3)$; the factor $x_2$ comes from the parameters, the locus
$x_2=0$ being the degenerate locus where $\rank H<\omega$, which is not part of the chart.

\textbf{What the fibre computation covers: the affine cone, not one chart.} The fibre over a direction
is a projective space $\mathbb P^{\tilde r}$, and the computation is \emph{not} performed on a single
affine chart of it. The parameter is a nonzero vector $x\in\mathbb R^{\tilde r}$, $\tilde r=r(v)$, that is,
the affine cone with the origin removed; projective points are the equivalence classes
$[x]=\{tx:\ t\ne0\}$. Each maximal minor $m_I(x,\lambda)$ is homogeneous of degree $\omega$ in $x$, so
the condition ``there is $\lambda$ with $m_I(x,\lambda)=0$ for all $I$'' is invariant under
$x\mapsto tx$ and defines a property of $[x]$ alone. The gcd is taken in $K(x_1,\dots,x_{\tilde r})[\lambda]$,
the rational function field of the cone; $\deg_\lambda\gcd=0$ there says that no common root exists at
a \emph{general} cone point, and the full-rank lemma upgrades this to \emph{every} cone point. Hence the
conclusion concerns the whole projective fibre, as claimed, and not one chart or one sample point.

\textbf{Three assertions, and which certificate supports which.} One must distinguish
(a) \emph{some point of the fibre is feasible}; (b) \emph{a dense open subset of it is feasible};
(c) \emph{the whole fibre is infeasible}. A single point with $\gcd=1$ establishes (a); openness of the
good set then gives (b); a positive-degree gcd together with the full-rank leading-coefficient lemma
establishes (c). The counterexample uses (a) for the direction $(1\!:\!1)$ and (c) for the five
non-eigen directions; the contrast between (a) and (c) \emph{is} the non-constancy of feasibility on
$E_1$.

\textbf{Why divisibility suffices for the whole fibre.} For a general family of polynomials a positive
degree of the gcd over the function field does \emph{not} imply that every specialization has a common
root, since the common factor may specialize to a nonzero constant. What licenses the conclusion is the
full-rank leading-coefficient lemma of Section~\ref{sec:decision}: at every point of the cone
$\rank H=\omega$, hence some maximal minor has $\lambda$-degree exactly $\omega$, so a nonconstant $g$
cannot specialize to a nonzero constant; every specialization therefore carries a nontrivial common
factor. Concretely, $g\mid m_I$ is the family of divisions
\begin{equation}\label{sec6:eq:division}
m_I=g\cdot q_I+r_I,\qquad \deg_\lambda r_I<\deg_\lambda g,
\end{equation}
all with $r_I=0$, verified for each maximal minor separately in exact rational arithmetic; the
quotients are of moderate size and the expanded identities are collected in Appendix~\ref{app:data}.
The computation is carried out in the polynomial ring $\mathbb{Q}(x_0,x_1,x_2)[\lambda]$ with
$\lambda$-degree at most $3$, so neither the gcd nor the divisions require any algebraic extension.

\textbf{The feasible direction.} For the eigen direction $[1\!:\!1]$, at the rational point
\begin{equation}\label{sec6:eq:x0}
\xi_0=(1,2,3,4)
\end{equation}
of the cone over the fibre $\mathbb{P}^{3}$ the greatest common divisor of all maximal minors is
$\gcd_I m_I(\xi_0,\lambda)=1$; Appendix~\ref{app:data} gives the specialized minors and an explicit
Bezout combination $\sum_I u_I(\lambda)\,m_I(\xi_0,\lambda)=1$. The chart construction of
Section~\ref{sec:decision} turns $\xi_0$ into a concrete row space, which is then a feasible design: it
satisfies both structural conditions and the spectral condition. A single point suffices for a good
direction by (a)$\Rightarrow$(b) above; for a bad direction the whole fibre must be covered, which is
(c).

\subsection{The precise wording of the conclusion}\label{sec6:sub:conclusion}

\begin{prop}[counterexample]\label{sec6:prop:nonconstant}
There is a rational instance with $n=8$, $p=2$ and $q=4$ in which the compression has a
two-dimensional real eigenspace $E_1$; inside $E_1$ the fibre over the direction $[1\!:\!1]$ contains
a feasible observer, whereas the fibres over the five non-eigen directions tested are all infeasible.
Hence feasibility is not constant on that eigenspace.
\end{prop}

\begin{cor}[failure of the representative-direction reduction]\label{sec6:cor:reduction}
The procedure ``take one eigenvector for each distinct real eigenvalue of the compression and decide
feasibility from that representative'' is not correct: the representative $(1\!:\!0)$ would declare the
instance of Proposition~\ref{sec6:prop:nonconstant} infeasible, whereas the instance is feasible.
\end{cor}

\textbf{Limitations.} The conclusion does \emph{not} prove that the repeated-eigenspace case is
computationally hard (it refutes a particular reduction, not the existence of a general algorithm); it
does \emph{not} show that geometric simplicity of the real spectrum is necessary for polynomial-time
decidability (that hypothesis controls the dimension of the direction parameter, and the example only
shows that this dimension cannot be ignored); and it asserts nothing about the untested directions of
$E_1$.

\textbf{Reproducibility.} The matrices $A$, $C$, $N$, $M$, $A_{21}$ and the characteristic polynomial
are given in full in Section~\ref{sec6:sub:instance}; the six decisions, the divisibility certificates
and the chart-point certificate are exact rational computations whose expanded forms, together with a
short reproduction recipe, are given in Appendix~\ref{app:data}.

\subsection{Relation to the fixed-parameter theorem}\label{sec6:sub:relation}

The counterexample shows that the direction parameter must retain the \emph{whole} projective eigenspace
$\mathbb{P}(E_\mu)$: the complexity theorem of Section~\ref{sec:complexity} takes geometric simplicity of
the real spectrum as a hypothesis, and the general decision problem for repeated eigenspaces lies
outside its scope; the direction is also a genuine degree of freedom which cannot be absorbed into the
row-space parameters.
